Những điểm yếu của textbook RSA xuất hiện vì hàm thô
\[
M \mapsto M^e \bmod N
\]
là tất định và có cấu trúc đại số rõ ràng. Để biến RSA thành một sơ đồ mã hóa thực tế, trước hết cần biến đổi thông điệp bằng một quá trình mã hóa tiền xử lý ngẫu nhiên. Quá trình này được gọi là \textbf{padding} (đệm).

\subsection{Vì sao cần đệm}

Padding phục vụ hai mục tiêu thiết yếu. Thứ nhất, nó đưa vào \textbf{tính ngẫu nhiên}, bảo đảm rằng cùng một thông điệp không luôn luôn cho ra cùng một bản mã. Thứ hai, nó thêm \textbf{cấu trúc}, để những đầu vào sai định dạng hoặc bị sửa đổi có thể được phát hiện và từ chối sau khi giải mã.

Về mặt khái niệm, đệm biến quá trình mã hóa từ
\[
C = M^e \bmod N
\]
thành
\[
C = \operatorname{RSA}\bigl(\operatorname{Pad}(M;r)\bigr) = \operatorname{Pad}(M;r)^e \bmod N,
\]
trong đó $r$ là ngẫu nhiên mới, được lấy độc lập cho mỗi lần mã hóa. Khi hàm RSA được áp dụng lên một biểu diễn đã ngẫu nhiên hóa của thông điệp thay vì lên chính thông điệp, tính tất định và cấu trúc nhân từng bị khai thác chống lại textbook RSA không còn dẫn tới các tấn công hữu dụng: kiểm tra bằng nhau, tấn công từ điển, khai căn với thông điệp nhỏ và thao tác bản mã đều trở nên khó hơn rất nhiều.

Một cách phát biểu mục tiêu thiết kế là padding phải biến RSA thành một sơ đồ \textbf{an toàn ngữ nghĩa} (semantic security), và lý tưởng là an toàn trước \textbf{tấn công bản mã được chọn thích nghi} (IND-CCA2), sao cho đối phương không học được gì từ một bản mã ngay cả khi được phép truy vấn một oracle giải mã trên các bản mã khác.

\subsection{PKCS\#1 v1.5}

Một sơ đồ đệm có ý nghĩa lịch sử quan trọng là \textbf{PKCS\#1 v1.5}. Gọi $k$ là độ dài tính bằng byte của mô-đun $N$. Trước khi mã hóa RSA, một thông điệp $M$ gồm $\mathrm{mLen}$ byte được nhúng vào một khối $k$ byte có dạng
\[
\underbrace{\mathtt{00}}_{1}\,\|\,\underbrace{\mathtt{02}}_{1}\,\|\,\underbrace{\text{PS}}_{\ge 8}\,\|\,\underbrace{\mathtt{00}}_{1}\,\|\,\underbrace{M}_{\mathrm{mLen}} .
\]
Các thành phần có vai trò cụ thể:

\begin{itemize}
\item byte \texttt{00} đứng đầu bảo đảm rằng khối, khi đọc như một số nguyên, nhỏ hơn $N$;
\item byte kiểu khối \texttt{02} cho biết khối này dành cho \emph{mã hóa};
\item chuỗi đệm PS gồm các byte ngẫu nhiên \textbf{khác 0} và phải dài ít nhất $8$ byte, đây chính là phần cung cấp tính ngẫu nhiên;
\item byte phân cách \texttt{00} đánh dấu nơi kết thúc phần đệm và bắt đầu thông điệp.
\end{itemize}

Vì PS phải chiếm ít nhất $8$ byte và các byte cố định chiếm thêm $3$ byte, độ dài thông điệp bị chặn bởi $\mathrm{mLen} \le k - 11$. Để giải mã, ta tính $C^d \bmod N$, biểu diễn lại thành khối $k$ byte, rồi kiểm tra tiền tố \texttt{00 02}, vùng đệm khác 0 và byte phân cách \texttt{00}; thông điệp là phần đứng sau đó.

Thiết kế này từng được triển khai rộng rãi và là một bước tiến lớn so với RSA giáo khoa. Tuy nhiên, về sau người ta chứng minh nó có thể bị tấn công bản mã được chọn thích nghi mỗi khi một cài đặt làm lộ việc khối sau giải mã có đúng định dạng PKCS\#1 hay không (xem Mục~7). Vì vậy, PKCS\#1 v1.5 minh họa một chủ đề lặp đi lặp lại trong mật mã: một quy tắc mã hóa trông mạnh trên giấy vẫn có thể thất bại khi kết hợp với một cài đặt làm rò rỉ dù chỉ một bit về giá trị sau giải mã.

\subsection{OAEP}

Một thiết kế hiện đại và có cơ sở lý thuyết hơn là \textbf{Optimal Asymmetric Encryption Padding (OAEP)}, do Bellare và Rogaway đề xuất. OAEP tiền xử lý thông điệp bằng ngẫu nhiên mới kết hợp với hai hàm sinh mặt nạ dựa trên hàm băm, ký hiệu $G$ và $H$, trước khi áp dụng phép lũy thừa RSA. Ở mức khái quát, nó:

\begin{itemize}
\item thêm một hạt giống ngẫu nhiên mới cho mỗi lần mã hóa,
\item trộn thông điệp và hạt giống qua một mạng kiểu Feistel hai vòng dựng từ $G$ và $H$, và
\item tạo ra một khối mã hóa có tính \emph{tất cả hoặc không gì}: chỉ có thể khôi phục thông điệp từ toàn bộ khối, không bao giờ từ một phần của nó.
\end{itemize}

Cụ thể, gọi $\mathrm{hLen}$ là độ dài đầu ra của hàm băm và $k$ là độ dài byte của $N$. Thông điệp trước hết được định dạng thành một khối dữ liệu
\[
\mathrm{DB} = \mathrm{lHash} \,\|\, \mathrm{PS} \,\|\, \mathtt{01} \,\|\, M,
\]
trong đó $\mathrm{lHash}$ là băm của một nhãn tùy chọn và $\mathrm{PS}$ là một dãy các byte 0. Sau đó lấy một hạt giống ngẫu nhiên dài $\mathrm{hLen}$, và thông điệp đã mã hóa được tính như sau:
\begin{align*}
\text{maskedDB} &= \mathrm{DB} \oplus G(\text{seed}), \\
\text{maskedSeed} &= \text{seed} \oplus H(\text{maskedDB}), \\
\mathrm{EM} &= \mathtt{00} \,\|\, \text{maskedSeed} \,\|\, \text{maskedDB},
\end{align*}
và $\mathrm{EM}$ là số nguyên được đưa vào phép hoán vị RSA. Giải mã chỉ việc đảo ngược hai vòng: $H$ khôi phục hạt giống từ maskedSeed và maskedDB, rồi $G$ khôi phục DB, sau đó kiểm tra định dạng của DB.

\begin{figure}[h]
\centering
\begin{tikzpicture}[>=Latex,node distance=0.6cm]
    \node[draw,rounded corners,fill=gray!10,minimum width=2.6cm,minimum height=0.9cm] (db)   at (0,4)   {DB ($\mathrm{lHash}\|\mathrm{PS}\|\mathtt{01}\|M$)};
    \node[draw,rounded corners,fill=gray!10,minimum width=2.0cm,minimum height=0.9cm] (seed) at (6.4,4) {seed (ngẫu nhiên)};

    \node[draw,circle,inner sep=1.5pt] (x1) at (0,2.4) {$\oplus$};
    \node[draw,rounded corners,fill=gray!10,minimum width=1.0cm,minimum height=0.8cm] (G) at (3.2,2.4) {$G$};

    \node[draw,rounded corners,fill=gray!10,minimum width=2.4cm,minimum height=0.9cm] (mdb) at (0,1.0) {maskedDB};

    \node[draw,rounded corners,fill=gray!10,minimum width=1.0cm,minimum height=0.8cm] (H) at (3.2,1.0) {$H$};
    \node[draw,circle,inner sep=1.5pt] (x2) at (6.4,1.0) {$\oplus$};

    \node[draw,rounded corners,fill=gray!10,minimum width=2.4cm,minimum height=0.9cm] (mseed) at (6.4,-0.4) {maskedSeed};

    \node[draw,rounded corners,fill=gray!12,minimum width=7.2cm,minimum height=0.9cm] (em) at (3.2,-1.9) {$\mathrm{EM}=\mathtt{00}\,\|\,\text{maskedSeed}\,\|\,\text{maskedDB}$};

    \draw[->,thick] (db.south) -- (x1.north);
    \draw[->,thick] (seed.south) .. controls (6.4,3.0) and (4.2,2.4) .. (G.east);
    \draw[->,thick] (G.west) -- (x1.east);
    \draw[->,thick] (x1.south) -- (mdb.north);
    \draw[->,thick] (mdb.east) -- (H.west);
    \draw[->,thick] (H.east) -- (x2.west);
    \draw[->,thick] (seed.east) .. controls (7.6,4) and (7.6,1.0) .. (x2.east);
    \draw[->,thick] (x2.south) -- (mseed.north);
    \draw[->,thick] (mdb.south) .. controls (0,-1.0) .. (em.north west);
    \draw[->,thick] (mseed.south) -- (em.north east);
\end{tikzpicture}
\caption{OAEP kết hợp khối dữ liệu DB và một hạt giống ngẫu nhiên qua hai hàm sinh mặt nạ $G$ và $H$ trước khi áp dụng RSA. Giải mã phải xử lý toàn bộ khối: muốn khôi phục $M$ cần toàn bộ maskedDB (để dựng lại hạt giống qua $H$) và toàn bộ hạt giống (để dựng lại DB qua $G$).}
\end{figure}

OAEP quan trọng vì nó đi kèm các bảo đảm lý thuyết mạnh hơn. Cấu trúc này có tính \emph{tất cả hoặc không gì} (all-or-nothing): lật một bit bất kỳ của $\mathrm{EM}$ sẽ làm xáo trộn không đoán trước được cả hạt giống lẫn khối dữ liệu khôi phục, nên kẻ tấn công không thể biến một bản mã hợp lệ thành một bản mã hợp lệ khác. Trong mô hình random oracle, RSA-OAEP được chứng minh là \textbf{an toàn IND-CCA2} dưới giả thuyết RSA, tức chính khả năng chống tấn công bản mã được chọn thích nghi mà PKCS\#1 v1.5 còn thiếu.

Bài học rộng hơn là \textbf{mã hóa RSA an toàn không chỉ đơn giản là ``RSA cộng một thông điệp''}. Nó là phép hoán vị RSA hợp thành với một thủ tục mã hóa ngẫu nhiên được thiết kế cẩn thận. Nếu thiếu bước tiền xử lý này, trapdoor permutation vẫn thanh nhã về toán học nhưng không an toàn về mật mã.
